\documentclass[10pt,aspectratio=169]{beamer}
% Preámbulo modular para presentaciones Beamer (Modelación Estadística)
% Diseño y paleta institucional del Tecnológico de Monterrey

\documentclass[aspectratio=169,xcolor={dvipsnames,table}]{beamer}

\usepackage[utf8]{inputenc}
\usepackage[spanish,mexico]{babel}
\usepackage{amsmath,amssymb,amsfonts,mathtools}
\usepackage{graphicx}
\usepackage{listings}
\usepackage{booktabs}
\usepackage{tikz}
\usetikzlibrary{matrix,arrows,decorations.pathmorphing,shapes.geometric,positioning}

% Paleta Institucional Tecnológico de Monterrey
\definecolor{TecAzul}{HTML}{0039A6}       % Azul TEC: color primario institucional
\definecolor{TecAzulOscuro}{HTML}{1A2E51} % Azul marino oscuro
\definecolor{TecRojo}{HTML}{EC2661}       % Rojo de acento (paleta miTec)
\definecolor{TecGris}{HTML}{646464}       % Gris neutro para comentarios y subtítulos
\definecolor{TecFondoBloque}{HTML}{F4F6F9}% Fondo suave para bloques

% Configuración de Tema Beamer
\usetheme{Boadilla}
\usecolortheme[named=TecAzulOscuro]{structure}
\setbeamercolor{alerted text}{fg=TecRojo}
\setbeamercolor{palette primary}{bg=TecAzulOscuro,fg=white}
\setbeamercolor{palette secondary}{bg=TecAzul,fg=white}
\setbeamercolor{palette tertiary}{bg=TecAzulOscuro!80!black,fg=white}
\setbeamercolor{title}{fg=TecAzulOscuro,bg=TecFondoBloque}
\setbeamercolor{frametitle}{fg=TecAzulOscuro,bg=TecFondoBloque}

% Bloques personalizados elegantes
\setbeamercolor{block title}{bg=TecAzulOscuro,fg=white}
\setbeamercolor{block body}{bg=TecFondoBloque,fg=black}
\setbeamercolor{block title alerted}{bg=TecRojo,fg=white}
\setbeamercolor{block body alerted}{bg=TecRojo!8,fg=black}
\setbeamercolor{block title example}{bg=TecAzul,fg=white}
\setbeamercolor{block body example}{bg=TecAzul!8,fg=black}

% Redondear bordes y añadir sombra sutil
\setbeamertemplate{blocks}[rounded][shadow=true]
\setbeamertemplate{navigation symbols}{}

% Pie de página limpio con número de diapositiva
\setbeamertemplate{footline}{
  \leavevmode%
  \hbox{%
  \begin{beamercolorbox}[wd=.7\paperwidth,ht=2.25ex,dp=1ex,left]{author in head/foot}%
    \hspace*{2ex}\insertshortauthor~~(\insertshortinstitute) \hfill \insertshorttitle
  \end{beamercolorbox}%
  \begin{beamercolorbox}[wd=.3\paperwidth,ht=2.25ex,dp=1ex,right]{date in head/foot}%
    \insertshortdate{}\hspace*{2em}
    \insertframenumber{} / \inserttotalframenumber\hspace*{2ex}
  \end{beamercolorbox}}%
  \vskip0pt%
}

% Configuración de Listings de Python para visualización pedagógica de código
\lstset{
    language=Python,
    basicstyle=\ttfamily\footnotesize,
    keywordstyle=\color{TecAzulOscuro}\bfseries,
    stringstyle=\color{TecRojo},
    commentstyle=\color{TecGris}\itshape,
    showstringspaces=false,
    breaklines=true,
    frame=lines,
    rulecolor=\color{TecAzulOscuro!30},
    backgroundcolor=\color{TecFondoBloque},
    aboveskip=0.5em,
    belowskip=0.5em,
    literate={á}{{\'a}}1 {é}{{\'e}}1 {í}{{\'i}}1 {ó}{{\'o}}1 {ú}{{\'u}}1
             {Á}{{\'A}}1 {É}{{\'E}}1 {Í}{{\'I}}1 {Ó}{{\'O}}1 {Ú}{{\'U}}1
             {ñ}{{\~n}}1 {Ñ}{{\~N}}1 {¿}{{?`}}1 {¡}{{!`}}1
}

% Metadatos institucionales por defecto
\title[Modelación Estadística]{Modelación Estadística para Ciencia de Datos e Ingeniería}
\author[J. Castillo Colmenares]{Juliho David Castillo Colmenares}
\institute[Tec de Monterrey]{Tecnológico de Monterrey \\ Escuela de Ingeniería y Ciencias}
\date{\today}
% Comandos y macros estadísticas compartidas para las presentaciones Beamer (Metropolis)

% Conjuntos numéricos básicos
\newcommand{\R}{\mathbb{R}}
\newcommand{\N}{\mathbb{N}}
\newcommand{\Q}{\mathbb{Q}}
\newcommand{\Z}{\mathbb{Z}}
\newcommand{\set}[1]{\left\{ #1 \right\}}
\newcommand{\abs}[1]{\left|#1\right|}
\newcommand{\norm}[1]{\left\| #1 \right\|}

% Comandos estadísticos y probabilísticos del curso
\newcommand{\Var}{\operatorname{Var}}
\newcommand{\cov}{\operatorname{Cov}}
\newcommand{\corr}{\rho}
\newcommand{\comb}[2]{\begin{pmatrix} #1 \\ #2 \end{pmatrix}}
\newcommand{\s}{\sigma}
\newcommand{\particion}{\mathfrak{P}}
\newcommand{\card}[1]{\#\left( #1 \right)}
\newcommand{\seq}[3]{\set{#1}_{#2}^{#3}}
\newcommand{\E}{\mathbb{E}}
\newcommand{\Prob}{\mathbb{P}}

% Entornos pedagógicos y cajas en español académico natural
\newenvironment{discusion}{%
  \begin{alertblock}{Pregunta de Discusión}%
}{%
  \end{alertblock}%
}

\newenvironment{reflexion}{%
  \begin{alertblock}{Reflexión para el Aula}%
}{%
  \end{alertblock}%
}

\newenvironment{paradoja}{%
  \begin{alertblock}{Paradoja de Exploración}%
}{%
  \end{alertblock}%
}

\newenvironment{motivacion}{%
  \begin{exampleblock}{Motivación: Ciencia de Datos en la Práctica}%
}{%
  \end{exampleblock}%
}

\newenvironment{retoclase}{%
  \begin{block}{Reto Rápido en Equipo}%
}{%
  \end{block}%
}

\title[Valores $p$ y decisiones]{Sección 07.03: Valores $p$ y decisiones}\subtitle{Evidencia, significación e interpretación}
\author[J. Castillo Colmenares]{Juliho Castillo Colmenares}\institute{Tecnológico de Monterrey}\date{}
\begin{document}
\begin{frame}[plain]\titlepage\end{frame}
\begin{frame}{Hoja de ruta}\small\begin{enumerate}\item Nivel de significación.\item Definición del valor $p$.\item Procedimiento e interpretación.\end{enumerate}\end{frame}
\begin{frame}{Fuente y alcance}\small La nota define el nivel $\alpha$, el valor $p$ y un procedimiento completo de decisión. Este mazo sigue ese recorrido.\begin{block}{Pregunta guía}¿Qué tan incompatibles son los datos con $H_0$?\end{block}\end{frame}
\begin{frame}{Nivel de significación}\small El nivel $\alpha$ es la probabilidad máxima tolerada de un error tipo I. Valores usuales son $0.10$, $0.05$ y $0.01$.\begin{alertblock}{Antes de observar} $\alpha$ se fija antes de calcular el estadístico.\end{alertblock}\end{frame}
\begin{frame}{Valor $p$}\small\begin{block}{Definición} El valor $p$ es la probabilidad de observar un resultado al menos tan extremo como el observado, suponiendo verdadera la hipótesis nula.\end{block}\end{frame}
\begin{frame}{Qué no significa}\small\begin{itemize}\item No es la probabilidad de que $H_0$ sea verdadera.\item No mide por sí solo el tamaño del efecto.\item Depende del modelo, la estadística y la dirección de la alternativa.\end{itemize}\end{frame}
\begin{frame}{Regla de decisión}\small\[p\leq\alpha\Rightarrow\text{rechazar }H_0;\qquad p>\alpha\Rightarrow\text{no rechazar }H_0.\]\begin{block}{Lenguaje} La decisión se refiere a la evidencia, no a una certeza lógica.\end{block}\end{frame}
\begin{frame}{Alternativas y colas}\small En una prueba unilateral, el valor $p$ se calcula en una cola; en una bilateral, se consideran ambos extremos. La alternativa determina la dirección.\end{frame}
\begin{frame}{Procedimiento I}\small\begin{enumerate}\item Formular $H_0$ y $H_a$.\item Elegir $\alpha$.\item Seleccionar la estadística y su distribución bajo $H_0$.\end{enumerate}\end{frame}
\begin{frame}{Procedimiento II}\small\begin{enumerate}\setcounter{enumi}{3}\item Calcular la estadística observada.\item Obtener el valor $p$.\item Comparar con $\alpha$.\end{enumerate}\end{frame}
\begin{frame}{Procedimiento III}\small\begin{block}{Reporte} Incluir estadística, grados de libertad, valor $p$, nivel $\alpha$ y conclusión contextual.\end{block}\end{frame}
\begin{frame}{Aplicación: planteamiento}\small Una empresa afirma que el tiempo medio de entrega es de tres días. Se sospecha que es mayor; se toma una muestra y se calcula un estadístico $t$.\begin{block}{Hipótesis} $H_0:\mu=3$ frente a $H_a:\mu>3$.\end{block}\end{frame}
\begin{frame}{Aplicación: estadística}\small Con $\bar x$, $s$ y $n$,\[t=\frac{\bar x-3}{s/\sqrt n}.\]Sus grados de libertad son $n-1$.\end{frame}
\begin{frame}{Aplicación: valor $p$}\small El valor $p$ es el área de la cola derecha más allá del estadístico observado bajo $t_{n-1}$.\end{frame}
\begin{frame}{Aplicación: decisión}\small Si $p\leq0.05$, se rechaza $H_0$ y se reporta evidencia de un tiempo medio mayor. Si $p>0.05$, no hay evidencia suficiente.\end{frame}
\begin{frame}{Interpretación I}\small La significación estadística no implica relevancia práctica. Conviene acompañar el valor $p$ con un intervalo y una medida del efecto.\end{frame}
\begin{frame}{Interpretación II}\small Un valor $p$ pequeño indica incompatibilidad con $H_0$ bajo el modelo; no demuestra que la alternativa sea verdadera ni que el resultado sea importante.\end{frame}
\begin{frame}{Consideraciones importantes}\small\begin{itemize}\item Elegir $\alpha$ antes del análisis.\item No cambiar la alternativa después de ver los datos.\item Evitar consultas repetidas sin control del error.\end{itemize}\end{frame}
\begin{frame}{Aplicación de reporte}\small\begin{block}{Plantilla} ``La estadística fue $T=\cdots$, con valor $p=\cdots$. Al nivel $\alpha=\cdots$, [rechazamos/no rechazamos] $H_0$.''\end{block}\end{frame}
\begin{frame}{Resumen del procedimiento}\small\begin{enumerate}\item Hipótesis.\item Nivel.\item Estadística.\item Valor $p$.\item Decisión.\item Interpretación.\end{enumerate}\end{frame}
\begin{frame}{Conclusión práctica}\small Una decisión reproducible conserva la alternativa, el nivel fijado y el valor $p$; así otra persona puede revisar tanto el cálculo como la interpretación.\end{frame}
\begin{frame}{Conexión con el capítulo}\small Los valores $p$ complementan los intervalos de confianza y preparan las pruebas para medias, proporciones y varianzas.\end{frame}
\end{document}
